\appendix
\section{Apêndice A: Derivação do IRD}

Considere uma amostra composta por $N$ indivíduos, classificados segundo $q$ características categóricas. Cada característica particiona o conjunto de indivíduos em categorias mutuamente exclusivas. Denote por $C$ o número total de categorias consideradas ao longo de todas as características.

Defina $X_{ic}$ como uma variável binária tal que $X_{ic} = 1$ se o indivíduo $i$ pertence à categoria $c$, e $X_{ic} = 0$ caso contrário. Por construção, cada indivíduo pertence exatamente a uma categoria por característica, de modo que:

\begin{equation}
    \sum_{c=1}^{C} X_{ic} = 1, \quad \forall i
\end{equation}

Sejam $p_c$ e $n_c$ as proporções da categoria $c$ na população e na amostra, respectivamente, satisfazendo:

\begin{equation}
    \sum_{c=1}^{C} p_c = \sum_{c=1}^{C} n_c = q
\end{equation}

O objetivo da ponderação é assegurar que, após a aplicação de pesos individuais 
$w_i$, as distribuições marginais da amostra ponderada coincidam com as proporções populacionais de interesse. Formalmente, impõe-se a condição de calibração:

\begin{equation}
    \frac{\sum_{i=1}^{N} w_i X_{ic}}{\sum_{i=1}^{N} w_i} = p_c, \quad \forall c = 1, \ldots, C
\end{equation}

Para simplificar a derivação, introduzem-se pesos não normalizados 
$\widetilde{w_i}$, de forma que a normalização posterior não altere as proporções relativas entre categorias. Assim, requer-se apenas que:

\begin{equation}
    \sum_{i=1}^{N} \widetilde{w_i} X_{ic} \propto p_c.
\end{equation}

\subsubsection*{Correção por uma característica}

Considere inicialmente o caso de uma única característica categórica. Supondo pesos constantes dentro de cada categoria, isto é, $\widetilde{w_i} = a_c$ para todo indivíduo pertencente à categoria $c$, obtém-se:
$\widetilde{w_i} = a_c$ para todo indivíduo pertencente à categoria $c$, obtém-se:

\begin{equation}
    \sum_{i=1}^{N} \widetilde{w_i} X_{ic} = a_c \sum_{i=1}^{N} X_{ic} = a_c N n_c.
\end{equation}

Para que essa quantidade seja proporcional a $p_c$, é necessário que:

\begin{equation}
    a_c N n_c \propto p_c.
\end{equation}

o que implica:

\begin{equation}
    a_c = \frac{p_c}{n_c}.
\end{equation}

Esse resultado estabelece que o fator de correção associado a cada categoria é dado pela razão entre sua proporção populacional e sua proporção observada na amostra.

\subsubsection*{Extensão para múltiplas características}

No caso geral, cada indivíduo pertence simultaneamente a $q$ categorias, uma para cada característica considerada. Assim, o peso individual deve refletir a contribuição de todas as categorias às quais o indivíduo pertence.

Uma forma linear e aditiva de combinar os fatores de correção é dada por:

\begin{equation}
    \widetilde{w_i} = \sum_{c=1}^{C} X_{ic} \frac{p_c}{n_c}.
\end{equation}


Essa especificação atribui a cada indivíduo a soma dos fatores de correção correspondentes às categorias que o caracterizam, preservando a contribuição independente de cada característica e evitando a amplificação excessiva dos pesos.

Observa-se que, dado que cada indivíduo contribui com exatamente $q$ termos na soma acima, a restrição $\sum_{c=1}^{C} p_c = \sum_{c=1}^{C} n_c = q$ garante a comparabilidade dos pesos entre indivíduos.

Após a normalização dos pesos $\widetilde{w_i}$, a condição de calibração é satisfeita, assegurando que as distribuições marginais da amostra ponderada reproduzam as proporções populacionais especificadas.
